\section*{Introduction}

The objective of this project is to dive into the practical application of \textbf{3D Slicer}, 
an open-source platform that has become a standard for medical image informatics and 
visualization. Rather than just following a tutorial, this work is designed as 
an \textbf{exploratory journey} through the process of transforming raw \textbf{DICOM data} 
into meaningful 3D representations.

Throughout the development of these tasks, we focused on learning the core concepts of 
volumetric processing, from the initial \textbf{segmentation} of anatomical structures 
to the creation of advanced \textbf{3D visualizations}. This report serves as a log of 
our progress, where we explain the methodology used to solve each exercise and, importantly, 
discuss the \textbf{troubles and challenges} we encountered along the way—such as refining 
complex geometries or navigating the vast ecosystem of \textbf{Slicer plugins}.

The project is structured into the following key areas:
\begin{itemize}
    \item \textbf{Segmentation and AI Exploration}: We compare manual segmentation techniques 
    with \textbf{AI-based modules} to understand the current capabilities and limitations of 
    automated medical tools.
    \item \textbf{Visualization and Metrics}: We explore how different \textbf{transfer functions} 
    and \textbf{markups} can be used to extract quantitative insights that would be useful for a 
    medical professional.
    \item \textbf{Plugin Integration}: We investigate advanced tools not covered in our standard 
    lab sessions to see how specialized plugins can tailor the software to specific clinical 
    tasks.
\end{itemize}

Ultimately, this project is about gaining hands-on experience and a deeper understanding of 
how to bridge the gap between raw medical data and \textbf{meaningful 3D insights}.